\section{Practical Implications for Code Translation}

The findings of this thesis suggest that modern code-specific large language
models can already provide useful support for real-world code migration tasks.
Both StarCoder2 and Qwen were able to generate syntactically valid and
partially functional translations for a considerable portion of the evaluation
dataset. Fine-tuning further improved performance, especially in terms of
accepted solutions and reduced compilation errors.

In practice, this means that such models can be used to accelerate software
migration projects by generating an initial translation that developers can
later refine. This may reduce the amount of manual effort required when
translating repetitive or relatively simple code fragments between programming
languages. For example, models may be particularly useful for translating small
utility functions, algorithmic code, or boilerplate structures. %szb: bizonyíték hiányzik

However, the results also show that current models are not yet reliable enough for
fully automated code migration. Even the best-performing model in this thesis still
produced a substantial number of compilation errors, wrong answers, and runtime
failures. Previous large-scale studies have similarly found that current LLM-based
code translation systems frequently generate incorrect translations, with correct
translation rates often remaining below 50\% depending on the model and language
pair \cite{llm_limitations}. Large-scale empirical studies further report that
LLM-based translation systems still struggle with complex programs, longer code
fragments, and semantic preservation, making human supervision necessary in
practical migration scenarios \cite{correctness_llm}.

Human review therefore remains necessary. Developers must still verify whether
the generated code compiles correctly, preserves the behavior of the original
program, and follows best practices of the target language. This is especially
important for larger codebases, where dependencies, project structure, and
runtime context introduce additional complexity.

Another important implication is that execution-based evaluation should be
integrated into practical migration workflows. Instead of relying only on
similarity-based metrics such as BLEU, generated translations should be tested
using compilation and execution-based methods whenever possible. This can help
identify semantic and logical errors that would otherwise remain hidden.

